
\documentclass{article}
\usepackage{colortbl}
\usepackage{makecell}
\usepackage{multirow}
\usepackage{supertabular}

\begin{document}

\newcounter{utterance}

\centering \large Interaction Transcript for game `cladder', experiment `full\_v1.5\_default', episode 2803 with Qwen3.5{-}9B.
\vspace{24pt}

{ \footnotesize  \setcounter{utterance}{1}
\setlength{\tabcolsep}{0pt}
\begin{supertabular}{c@{$\;$}|p{.15\linewidth}@{}p{.15\linewidth}p{.15\linewidth}p{.15\linewidth}p{.15\linewidth}p{.15\linewidth}}
    \# & \multicolumn{2}{c}{Player} && \multicolumn{2}{c}{Game Master} \\
    \hline

    \theutterance \stepcounter{utterance}  
    & & & \multicolumn{4}{p{0.6\linewidth}}{
        \cellcolor[rgb]{0.9,0.9,0.9}{
            \makecell[{{p{\linewidth}}}]{
                \texttt{\tiny{[P1$\langle$GM]}}
                \texttt{You are an expert in causal inference. The following question is not a typical commonsense query, but rather a meticulously designed question created by a professor specializing in causal inference, intended to assess the students' mastery of the course content.} \\
\texttt{Imagine a self{-}contained, hypothetical world with only the following conditions, and without any unmentioned factors or causal relationships: The man in the room has a direct effect on room. The candle has a direct effect on room. We know that blowing out the candle and candle with wax causes dark room. We observed the candle is out of wax. Would the room is bright if blowing out the candle instead of not blowing out the candle?} \\
\texttt{Start your answer with 'Yes' or 'No', followed by additional reasoning or evidence to support your explanation.} \\
            }
        }
    }
     \\ \\

    \theutterance \stepcounter{utterance}  
    & \multicolumn{4}{p{0.6\linewidth}}{
        \cellcolor[rgb]{0.561,0.737,0.561}{
            \makecell[{{p{\linewidth}}}]{
                \texttt{\tiny{[P1$\rangle$GM]}}
                \texttt{No.} \\
\\ 
\texttt{**Reasoning:**} \\
\\ 
\texttt{Based on the causal structure defined in your hypothetical world, we can deduce the following:} \\
\\ 
\texttt{1.  **Causal Mechanism**: The prompt explicitly states that "blowing out the candle and candle with wax causes dark room." This establishes a direct causal link where the state of the candle (specifically, being blown out while having wax) is a sufficient cause for the room being dark.} \\
\texttt{2.  **Observed State**: We are told the "candle is out of wax." In causal inference terms, this is the observed state of the variable \$C\$ (Candle).} \\
\texttt{3.  **Counterfactual Query**: The question asks for the outcome in the counterfactual scenario: "Would the room be bright if blowing out the candle instead of not blowing out the candle?"} \\
\texttt{    *   The condition "blowing out the candle" combined with the existing fact that the candle is "out of wax" (or simply the act of blowing it out in this specific world where that action triggers darkness) activates the cause identified in step 1.} \\
\texttt{    *   Since the cause ("blowing out the candle" + "candle with wax") leads to the effect ("dark room"), the resulting state of the room must be **dark**.} \\
\texttt{4.  **Conclusion**: The question asks if the room would be **bright**. Since the causal mechanism dictates the room becomes **dark** under the specified conditions, the answer is negative. The presence of the candle (even if out of wax) and the act of blowing it out in this specific world definition results in darkness, not brightness. The "direct effect on room" mentioned for the man and candle implies these are the primary drivers, and the specific rule provided overrides any assumption of brightness.} \\
\\ 
\texttt{Therefore, the room would not be bright; it would be dark.} \\
            }
        }
    }
    & & \\ \\

    \theutterance \stepcounter{utterance}  
    & & & \multicolumn{2}{p{0.3\linewidth}}{
        \cellcolor[rgb]{0.9,0.9,0.9}{
            \makecell[{{p{\linewidth}}}]{
                \texttt{\tiny{[GM$|$GM]}}
                \texttt{No} \\
            }
        }
    }
    & & \\ \\

    \theutterance \stepcounter{utterance}  
    & & & \multicolumn{2}{p{0.3\linewidth}}{
        \cellcolor[rgb]{0.9,0.9,0.9}{
            \makecell[{{p{\linewidth}}}]{
                \texttt{\tiny{[GM$|$GM]}}
                \texttt{yes} \\
            }
        }
    }
    & & \\ \\

    \theutterance \stepcounter{utterance}  
    & & & \multicolumn{2}{p{0.3\linewidth}}{
        \cellcolor[rgb]{0.9,0.9,0.9}{
            \makecell[{{p{\linewidth}}}]{
                \texttt{\tiny{[GM$|$GM]}}
                \texttt{game\_result = LOSE} \\
            }
        }
    }
    & & \\ \\

\end{supertabular}
}

\end{document}
